% Compiled with LuaLaTeX
\documentclass[tate, book, paper=a5]{jlreq}
\usepackage{luatexja-fontspec}
\usepackage{tikz}
\usepackage{eso-pic}
\usepackage{geometry}

% Margin settings
\geometry{margin=20mm}

% Subtle background design using eso-pic and TikZ
\AddToShipoutPictureBG{
    \begin{tikzpicture}[remember picture, overlay]
        % Define a very subtle color
        \colorlet{subtlecolor}{gray!30!white}
        
        % Top Right Corner (Subtle Meander/雷紋 motif)
        \draw[line width=0.4pt, color=subtlecolor] 
            ([shift={(-15mm, -15mm)}]current page.north east) 
            -- ++(-5mm, 0) -- ++(0, -5mm) -- ++(3mm, 0) -- ++(0, 3mm) -- ++(-1mm, 0);

        % Bottom Left Corner
        \draw[line width=0.4pt, color=subtlecolor] 
            ([shift={(15mm, 15mm)}]current page.south west) 
            -- ++(5mm, 0) -- ++(0, 5mm) -- ++(-3mm, 0) -- ++(0, -3mm) -- ++(1mm, 0);
    \end{tikzpicture}
}

\begin{document}

\chapter*{走れメロス}

メロスは激怒した。必ず、かの邪智暴虐の王を除かなければならぬと決意した。メロスには政治がわからぬ。メロスは、村の牧人である。笛を吹き、羊と遊んで暮して来た。けれども邪悪に対しては、人一倍に敏感であった。

きょう未明メロスは村を出発し、野を越え山越え、十里はなれた此のシラクスの市にやって来た。メロスには父も、母も無い。女房も無い。十六の、内気な妹と二人暮しだ。この妹は、村の或る律気な一牧人を、近々、花婿として迎える事になっていた。結婚式も間近かなのである。メロスは、それゆえ、花嫁の衣裳やら祝宴の御馳走やらを買いに、はるばる市にやって来たのだ。先ず、その品々を買い集め、それから都の大路をぶらぶら歩いた。メロスには竹馬の友があった。セリヌンティウスである。今は此のシラクスの市で、石工をしている。その友を、これから訪ねてみるつもりなのだ。久しく逢わなかったのだから、訪ねて行くのが楽しみである。

\end{document}